\section*{Bab \ref{ch:b3:environmental-toxicology} --- Kimia Lingkungan dan Toksikologi}

\begin{solution}{exo:b3:environmental-toxicology:1}
Molekul diatomik homonuklir memiliki satu vibrasi tunggal yang mempertahankan momen dipol nol,
dan argon tidak memiliki vibrasi: tidak satu pun dapat menyerap radiasi inframerah. $\ce{CO2}$ menyerap melalui
modus tekuknya (dan ulur antisimetrisnya), yang menimbulkan dipol yang
berosilasi; tekukan itu terletak di tengah-tengah emisi Bumi.
\end{solution}

\begin{solution}{exo:b3:environmental-toxicology:2}
$10^{0.1} = 1.26$: $[\ce{H+}]$ naik 26\,\%.
\end{solution}

\begin{solution}{exo:b3:environmental-toxicology:3}
$(2.0 + 0.50)\ \unit{mmol/L} \times \qty{100.1}{mg/mmol} = \qty{250}{mg/L}$ $\ce{CaCO3}$.
\end{solution}

\begin{solution}{exo:b3:environmental-toxicology:4}
N: $0.60/14.0 = \qty{0.043}{mmol/L}$; P: $0.010/31.0 = \qty{3.2e-4}{mmol/L}$; N\,:\,P $= 133$, jauh di atas 16: fosfor habis
lebih dahulu dan membatasi pertumbuhan.
\end{solution}

\begin{solution}{exo:b3:environmental-toxicology:5}
$L_0 = \mathrm{BOD_5}/(1 - \eu^{-0.23 \times 5}) = 170/0.683 = \qty{250}{mg/L}$.
\end{solution}

\begin{solution}{exo:b3:environmental-toxicology:6}
$\qty{4.40e-3}{L} \times \qty{0.0200}{mol/L} = \qty{8.80e-5}{mol}$ dalam \qty{0.1000}{L}: \omterm{def:b3:environmental-toxicology:alkalinity}{alkalinitas} $\qty{8.80e-4}{mol/L}$; sebagai $\ce{HCO3-}$ (\qty{61.0}{g/mol}),
\qty{53.7}{mg/L}.
\end{solution}

\begin{solution}{exo:b3:environmental-toxicology:7}
$K_d = 0.020 \times 300 = \qty{6.0}{L/kg}$. Tersorpsi/total $= K_dm_s/(K_dm_s + V_w) = 9.0/(9.0 + 0.30) = 0.97$: 97\,\% tersorpsi.
\end{solution}

\begin{solution}{exo:b3:environmental-toxicology:8}
Benzena: $\log\mathrm{BCF} = 0.85 \times 2.13 - 0.70 = 1.11$, BCF sekitar 13. PCB: $0.85 \times 6.67 - 0.70 = 4.97$, BCF sekitar
\num{9e4}: PCB itu memekat dalam ikan sekitar tujuh ribu kali lebih banyak daripada benzena; pada
$K_{ow}$ setinggi itu, regresinya mendekati batas keberlakuannya.
\end{solution}

\begin{solution}{exo:b3:environmental-toxicology:9}
$\qty{192}{mg/kg} \times \qty{70}{kg} = \qty{13}{g}$, sekitar 150 cangkir. Spesies-spesies berbeda dalam penyerapan dan metabolisme,
dan $\mathrm{LD_{50}}$ bukanlah ambang yang aman; efek merugikan sudah dimulai jauh di bawahnya.
\end{solution}

\begin{solution}{exo:b3:environmental-toxicology:10}
PNEC $= \qty{0.50}{mg/L}/1000 = \qty{0.50}{\micro g/L}$; $\mathrm{RQ} = 0.20/0.50 = 0.40$, di bawah 1: tidak ada kekhawatiran pada paparan ini.
\end{solution}

\begin{solution}{exo:b3:environmental-toxicology:11}
Plankton/air: $\qty{0.015}{mg/kg}/\qty{2.0e-6}{mg/L} = \qty{7.5e3}{L/kg}$ (biokonsentrasi). Ikan/plankton: $0.30/0.015 = 20$;
burung/ikan: $6.0/0.30 = 20$ (\omterm{def:b3:environmental-toxicology:bio}{biomagnifikasi} pada setiap tahap); burung/air: $\num{3e6}$.
\end{solution}

\begin{solution}{exo:b3:environmental-toxicology:12}
$\qty{1.0e9}{g}/\qty{64.1}{g/mol} = \qty{1.56e7}{mol}$ $\ce{SO2}$, yang memberikan jumlah mol $\ce{H2SO4}$ yang sama: $\qty{1.53e9}{g}$, sekitar \qty{1500}{t}. Asam itu
melepaskan $\qty{3.12e7}{mol}$ $\ce{H+}$; pada pH 4.0, $[\ce{H+}] = \qty{1.0e-4}{mol/L}$: $\qty{3.1e11}{L}$ hujan.
\end{solution}

\begin{solution}{pb:b3:environmental-toxicology:1}
\textbf{1.} Pestisida itu sepuluh ribu kali lebih larut dalam oktanol daripada dalam air: karena hidrofobik,
pestisida itu akan tersorpsi pada bahan organik dan menumpuk dalam lemak.

\textbf{2.} $K_{oc} = 0.41 \times 10^4 = \qty{4100}{L/kg}$; $K_{sw} = f_{oc}K_{oc} = 0.040 \times 4100 = \qty{164}{L/kg}$.

\textbf{3.} $K_{fw} = 0.048 \times 10^4 = \qty{480}{L/kg}$.

\textbf{4.} Tidak: pada kesetimbangan, udara menampung $10^{-4}$ kali konsentrasi dalam air.

\textbf{5.} Molekul netral yang hidrofobik itu larut dalam bahan organik, seperti dalam
oktanol; permukaan mineral bersifat polar dan tertutup air.

\textbf{6.} Sedimen, yang kapasitasnya, $K_{sw}m_s$, paling besar.

\textbf{7.} $V_w = \qty{1.0e9}{L}$; $K_{aw}V_a = \qty{1.0e8}{L}$; $K_{sw}m_s = \qty{1.64e10}{L}$; $K_{fw}m_f = \qty{4.8e6}{L}$.

\textbf{8.} $C_w = \qty{1.0e8}{mg}/\qty{1.75e10}{L} = \qty{5.7e-3}{mg/L}$, \qty{5.7}{\micro g/L}.

\textbf{9.} Air \qty{5.7}{kg}, udara \qty{0.57}{kg}, sedimen \qty{94}{kg}, ikan \qty{0.027}{kg}.

\textbf{10.} Udara $\qty{5.7e-7}{mg/L}$; sedimen \qty{0.94}{mg/kg}; ikan \qty{2.7}{mg/kg}.

\textbf{11.} Kesetimbangan antara semua kompartemen sekaligus, tanpa penguraian, tanpa
aliran masuk atau keluar, dan dengan kompartemen-kompartemen yang tercampur sempurna.

\textbf{12.} Dengan penguraian dan aliran pada keadaan tunak (tingkat II), inventarisnya ditentukan oleh
laju masukan dan kehilangan; dengan laju transfer antarkompartemen (tingkat III),
konsentrasi-konsentrasinya tidak lagi berada pada rasio kesetimbangan, dan kompartemen yang
menerima masukan menjadi lebih kaya.

\textbf{13.} $k = \ln 2/\qty{60}{d} = \qty{0.0116}{d^{-1}}$.

\textbf{14.} $\eu^{-0.0116 \times 365} = 0.0145$: tersisa 1.5\,\%.

\textbf{15.} $\qty{5.7}{\micro g/L} \times 0.0145 = \qty{0.083}{\micro g/L}$.

\textbf{16.} Penguraian sering lebih lambat di sedimen yang dingin dan anoksik, dan pestisida
yang tersorpsi dilepaskan kembali ke air secara perlahan, sehingga tetap ada setelah
airnya bersih.

\textbf{17.} Ikatan yang tahan terhadap hidrolisis, oksidasi, dan serangan mikroba (misalnya C--Cl
aromatik), kelarutan dalam air yang rendah, dan sorpsi yang melindunginya dari
pengurai.

\textbf{18.} PNEC $= \qty{50}{\micro g/L}/10 = \qty{5.0}{\micro g/L}$.

\textbf{19.} $\mathrm{RQ} = 5.7/5.0 = 1.1$.

\textbf{20.} Sedikit di atas 1: ada indikasi risiko, dan penilaiannya harus diperbaiki (data
toksisitas yang lebih baik, konsentrasi terukur) atau paparannya dikurangi.

\textbf{21.} \qty{2.7}{mg/kg}, hampir tiga kali ambang \qty{1.0}{mg/kg}: ikan itu tidak boleh dimakan
sampai konsentrasinya turun.

\textbf{22.} Konsentrasi dalam air, sedimen, dan ikan dari waktu ke waktu dan di beberapa titik,
untuk memeriksa model dan mengikuti penurunannya.

\textbf{23.} Lebih sedikit pestisida yang diaplikasikan, aplikasi yang dijauhkan dari hujan dan dari tepi danau,
jalur penyangga yang menahan limpasan, atau alternatif yang kurang persisten.

\textbf{24.} Faktor itu membagi tingkat tanpa efek yang terukur untuk menutup kesenjangan antara beberapa
spesies laboratorium dan seluruh ekosistem; faktor 10 alih-alih 1000 mengubah
PNEC, dan putusannya, seratus kali lipat.

\textbf{25.} \omterm{def:b3:environmental-toxicology:ecotoxicology}{Hasil bagi risiko} PEC/PNEC sekitar 1.1, sedikit di atas 1.
\end{solution}
